\chapter{Wave Equation}


\section{One-dimensional wave equation}

We consider
$
u_{tt}-\Delta u=0,
$
and more generally
$
u_{tt}-\Delta u=f
\text{ (inhomogeneous)}.
$

\begin{definition}
For
$
u:\overline{U}\times [0,\infty)\to \mathbb{R},
$
we denote the differential operator
$$
\square u := u_{tt}-\Delta u.
$$
\end{definition}

\begin{theorem}[Solution of the wave equation when $n=1$]
Assume
$
g\in C^2(\mathbb{R}), \qquad h\in C^1(\mathbb{R}).
$
Consider the initial value problem
$$
\begin{cases}
u_{tt}-u_{xx}=0 & \text{in } \mathbb{R}\times (0,\infty),\\
u=g & \text{on } \mathbb{R}\times \{t=0\},\\
u_t=h & \text{on } \mathbb{R}\times \{t=0\}.
\end{cases}
$$
Define $u$ by d'Alembert's formula:
$$
u(x,t)
=
\frac12\bigl[g(x+t)+g(x-t)\bigr]
+
\frac12\int_{x-t}^{x+t} h(y)\,dy,
\qquad (x\in \mathbb{R},\ t>0).
$$

Then
\begin{enumerate}
    \item $u\in C^2(\mathbb{R}\times [0,\infty))$,
    \item $u$ solves the IVP of the wave equation,
    \item
    $
    \lim_{(x,t)\to (x^0,0)} u(x,t)=g(x^0),
    \quad
    \lim_{(x,t)\to (x^0,0)} u_t(x,t)=h(x^0).
    $
\end{enumerate}
\end{theorem}

\begin{remark}
Factor the differential operator:
$$
\square
=
\partial_t^2-\partial_x^2
=
(\partial_t+\partial_x)(\partial_t-\partial_x).
$$
So one can solve a coupled system of two first-order equations:
$$
\begin{cases}
(\partial_t-\partial_x)u=v,\\
(\partial_t+\partial_x)v=0.
\end{cases}
$$
\end{remark}

\section{Higher dimensions: spherical means}

Now move to higher dimensions, first consider $n=3$.

\begin{definition}
Let $x\in \mathbb{R}^n$, $t>0$, $r>0$.
We denote
$$
U(x,r,t):=\fint_{\partial B(x,r)} u(y,t)\,dS(y).
$$
Similarly, define
$$
G(x,r):=\fint_{\partial B(x,r)} g(y)\,dS(y),
$$
and
$$
H(x,r):=\fint_{\partial B(x,r)} h(y)\,dS(y).
$$
\end{definition}

\begin{lemma}[Euler--Poisson--Darboux equation]
Fix $x\in \mathbb{R}^n$, and let $u$ satisfy the IVP for the $n$-dimensional wave equation.
Then for this fixed $x$, $U$ is a function of $r,t$, and we have
$$
U\in C^m(\overline{\mathbb{R}_+}\times [0,\infty))
$$
and $U$ satisfies
$$
\begin{cases}
U_{tt}-U_{rr}-\dfrac{n-1}{r}U_r=0 & \text{in } \mathbb{R}_+\times (0,\infty),\\[1em]
U=G & \text{on } \mathbb{R}_+\times \{t=0\},\\
U_t=H & \text{on } \mathbb{R}_+\times \{t=0\}.
\end{cases}
$$
\end{lemma}

\begin{lemma}[Reduction when $n=3$]
For $n=3$, set
$$
\tilde{U}=rU, \qquad \tilde{G}=rG, \qquad \tilde{H}=rH.
$$
Then $\tilde{U}$ solves
$$
\begin{cases}
\tilde{U}_{tt}-\tilde{U}_{rr}=0 & \text{in } \mathbb{R}_+\times (0,\infty),\\
\tilde{U}=\tilde{G} & \text{on } \mathbb{R}_+\times \{t=0\},\\
\tilde{U}_t=\tilde{H} & \text{on } \mathbb{R}_+\times \{t=0\},\\
\tilde{U}=0 & \text{on } \{r=0\}\times (0,\infty).
\end{cases}
$$
Also, $\tilde{U}$ can be expressed by
$$
\tilde{U}(x;r,t)
=
\frac12\bigl[\tilde{G}(r+t)-\tilde{G}(t-r)\bigr]
+
\frac12\int_{-r+t}^{\,r+t}\tilde{H}(y)\,dy.
$$
\end{lemma}

\begin{remark}
This can be solved by the odd reflection method.
\end{remark}


\begin{theorem}[Kirchhoff's formula]
For the $n=3$ IVP for the wave equation,
$$
u(x,t)
=
\frac{\partial}{\partial t}
\left(
t\fint_{\partial B(x,t)} g\,dS
\right)
+
t\fint_{\partial B(x,t)} h\,dS.
$$
Equivalently,
$$
u(x,t)
=
\fint_{\partial B(x,t)}
\bigl(
t\,h(y)+g(y)+Dg(y)\cdot (y-x)
\bigr)\,dS(y).
$$
\end{theorem}


\begin{remark}
\begin{enumerate}\
    \item At $(x,t)$, the solution depends only on $g,h$ evaluated on
    $
    \partial B(x,t)\subset \mathbb{R}^3.
    $
    \item The same approach cannot go through on $\mathbb{R}^2$.
    Instead, we will construct solutions on $\mathbb{R}^2$ by first finding an
    ``$x_3$-free'' solution on $\mathbb{R}^3$ and then restricting to $\mathbb{R}^2$.
\end{enumerate}
\end{remark}


\begin{theorem}[Poisson formula]
If $u$ is a smooth solution to
$$
\begin{cases}
u_{tt}-\Delta u=0 & \text{in } \mathbb{R}^2\times (0,\infty),\\
u|_{t=0}=g & \text{on } \mathbb{R}^2,\\
u_t|_{t=0}=h & \text{on } \mathbb{R}^2,
\end{cases}
$$
then
$$
u(x,t)
=
\frac12
\int_{B(x,t)}
\frac{t\,g(y)+t^2 h(y)+t\,Dg(y)\cdot (y-x)}
{\bigl(t^2-|y-x|^2\bigr)^{1/2}}
\,dy.
$$
\end{theorem}

\begin{proof}
Hadamard's method of descent.
\end{proof}

\begin{remark}
These ideas generalize to higher dimensions; see reading.
\end{remark}

\begin{theorem}[Solution of the nonhomogeneous wave equation]
Consider the inhomogeneous problem
$$
\begin{cases}
u_{tt}-\Delta u=f, & t>0,\\
u|_{t=0}=u_t|_{t=0}=0, & t=0.
\end{cases}
$$
Motivated by Duhamel's principle, define $\widetilde{u}=\widetilde{u}(x,t;s)$ to be the solution of
$$
\begin{cases}
\widetilde{u}_{tt}(\cdot;s)-\Delta \widetilde{u}(\cdot;s)=0
& \text{in } \mathbb{R}^n\times (s,\infty),\\
\widetilde{u}(\cdot;s)=0
& \text{on } \mathbb{R}^n\times \{t=s\},\\
\widetilde{u}_t(\cdot;s)=f(\cdot,s)
& \text{on } \mathbb{R}^n\times \{t=s\}.
\end{cases}
$$
Set
$$
u(x,t):=\int_0^t \widetilde{u}(x,t;s)\,ds.
$$
Then
\begin{enumerate}
    \item $u\in C^2(\mathbb{R}^n\times [0,\infty))$,
    \item $u$ is the solution of
    $
    u_{tt}-\Delta u=f \text{ in } \mathbb{R}^n\times (0,\infty),
    $
    \item
    $
    \lim_{(x,t)\to (x^0,0)}u(x,t)=0,
    \quad
    \lim_{(x,t)\to (x^0,0)}u_t(x,t)=0.
    $
\end{enumerate}
\end{theorem}

\section{Fourier method}

\begin{definition}[Fourier transform and inverse Fourier transform]
For $u\in L^1(\mathbb{R}^n)$, $n\ge 1$, define
$$
\widehat{u}(\xi)
=
\frac{1}{(2\pi)^{n/2}}
\int_{\mathbb{R}^n}
e^{-ix\cdot \xi}u(x)\,dx.
$$
This is the Fourier transform; denote $\mathcal{F}u=\widehat{u}$.
Define also
$$
\check{u}(x)
=
\frac{1}{(2\pi)^{n/2}}
\int_{\mathbb{R}^n}
e^{ix\cdot \xi}u(\xi)\,d\xi.
$$
This is the inverse Fourier transform; denote $\mathcal{F}^{-1}u=\check{u}$.
\end{definition}

\begin{theorem}[Plancherel's theorem]
Assume
$
u\in L^1(\mathbb{R}^n)\cap L^2(\mathbb{R}^n).
$
Then
$
\widehat{u},\check{u}\in L^2(\mathbb{R}^n),
$
and
$
\|\widehat{u}\|_{L^2(\mathbb{R}^n)}
=
\|\check{u}\|_{L^2(\mathbb{R}^n)}
=
\|u\|_{L^2(\mathbb{R}^n)}.
$
\end{theorem}

\begin{remark}\
\begin{itemize}
    \item The Fourier transform preserves the $L^2$ norm.
    \item This allows one to define the Fourier transform on $L^2$ functions by approximating
    $
    u\in L^2
    $
    by a sequence
    $
    (u_n)\subset L^1\cap L^2.
    $
    \item In fact, the Fourier transform can be defined on the Schwartz space.
\end{itemize}
\end{remark}



\begin{theorem}[Properties of the Fourier transform]
Assume $u,v\in L^1(\mathbb{R}^n)$.
Then:
\begin{enumerate}
    \item
    $
    \int_{\mathbb{R}^n} u\,\overline{v}\,dx
    =
    \int_{\mathbb{R}^n} \widehat{u}\,\overline{\widehat{v}}\,d\xi.
    $
    \item
    $
    (D^\alpha u)\,\widehat{\ } = (i\xi)^\alpha \widehat{u}.
    $
    \item If $u,v\in L^1\cap L^2$, then
    $
    (u*v)\,\widehat{\ } = (2\pi)^{n/2}\widehat{u}\,\widehat{v}.
    $
    \item
    $
    u=(\widehat{u})^\vee.
    $
\end{enumerate}
\end{theorem}

\begin{theorem}[Solution of the homogeneous wave equation by Fourier transform]
Consider the homogeneous IBVP
$$
\begin{cases}
u_{tt}-\Delta u=0 & \text{in } \mathbb{R}^n\times (0,\infty),\\
u|_{t=0}=g & \text{on } \mathbb{R}^n\times \{t=0\},\\
u_t|_{t=0}=h & \text{on } \mathbb{R}^n\times \{t=0\}.
\end{cases}
$$
Then
$$
u(t,x)
=
\cos(t|\nabla|)\,g
+
\frac{\sin(t|\nabla|)}{|\nabla|}\,h,
$$
where
$$
\cos(t|\nabla|)\,g
:=
\bigl(\cos(t|\xi|)\widehat{g}(\xi)\bigr)^\vee,
$$
and
$$
\frac{\sin(t|\nabla|)}{|\nabla|}\,h
:=
\left(
\frac{\sin(t|\xi|)}{|\xi|}\widehat{h}(\xi)
\right)^\vee.
$$
\end{theorem}

\begin{proof}
Let us do the Fourier transform in the $x$ variable:
$$
u(t,x)\mapsto \widehat{u}(t,\xi).
$$
Then
$$
\begin{cases}
\widehat{u}_{tt}+|\xi|^2\widehat{u}=0,\\
\widehat{u}|_{t=0}=\widehat{g},\\
\widehat{u}_t|_{t=0}=\widehat{h}.
\end{cases}
$$
For each fixed $\xi$, this leads to an ODE for
$
\gamma(t)=\widehat{u}(t,\xi):
$
$$
\begin{cases}
\ddot{\gamma}+|\xi|^2\gamma=0,\\
\gamma(0)=\widehat{g}(\xi),\\ 
\dot{\gamma}(0)=\widehat{h}(\xi).
\end{cases}
$$
This has solution
$$
\gamma(t)
=
\widehat{g}(\xi)\cos(t|\xi|)
+
\widehat{h}(\xi)\frac{\sin(t|\xi|)}{|\xi|}.
$$
Hence
$$
\widehat{u}(t,\xi)
=
\widehat{g}(\xi)\cos(t|\xi|)
+
\widehat{h}(\xi)\frac{\sin(t|\xi|)}{|\xi|}.
$$
Therefore
$$
u(t,x)=\bigl(\widehat{u}(t,\xi)\bigr)^\vee.
$$
\end{proof}

\begin{theorem}[Solution of the inhomogeneous wave equation by Fourier transform]
Consider the inhomogeneous IBVP
$$
\begin{cases}
u_{tt}-\Delta u=f & \text{on } \mathbb{R}^n\times (0,\infty),\\
u|_{t=0}=g,\qquad u_t|_{t=0}=h
& \text{in } \mathbb{R}^n\times \{t=0\}.
\end{cases}
$$
Then the problem is solved by
$$
u(t,x)
=
\cos(t|\nabla|)\,g
+
\frac{\sin(t|\nabla|)}{|\nabla|}\,h
+
\int_0^t
\frac{\sin((t-s)|\nabla|)}{|\nabla|}f(s)\,ds.
$$
\end{theorem}

\section{Energy method: uniqueness and finite propagation speed}

\begin{theorem}[Uniqueness for the wave equation]
Let $U\subset \mathbb{R}^n$ be bounded and open, with smooth boundary $\partial U$.
Consider the following IBVP:
$$
\begin{cases}
u_{tt}-\Delta u=f & \text{in } U_T,\\
u=g & \text{on } \Gamma_T,\\
u_t=h & \text{on } U\times \{t=0\}.
\end{cases}
$$
Then there exists at most one solution
$
u\in C^2(\overline{U_T}).
$
\end{theorem}

\begin{theorem}[Finite propagation speed]
Suppose
$
u\in C^2(\mathbb{R}^n)
$
solves
$$
\begin{cases}
u_{tt}-\Delta u=0 & \text{in } \mathbb{R}^n\times (0,\infty),\\
u(x,0)=u_t(x,0)=0 & \text{on } B(x_0,t_0)\times \{t=0\}.
\end{cases}
$$
Then
$
u\equiv 0
$
within the cone
$$
K(x_0,t_0)
=
\left\{
(x,t): 0\le t\le t_0,\ |x-x_0|\le t_0-t
\right\}.
$$
\end{theorem}